\documentclass[a4paper,12pt]{article}
\usepackage[utf8]{inputenc}
\usepackage[T1]{fontenc}
\usepackage[french]{babel}
\usepackage{geometry}
\geometry{left=2.5cm,right=2.5cm,top=2.5cm,bottom=2.5cm}

\usepackage{lmodern}
\usepackage{xcolor}
\definecolor{titleblue}{RGB}{0,51,140}

\usepackage{hyperref}

\pagestyle{empty}

\begin{document}

\begin{flushleft}
    \textbf{Ismail AISSAOUI} \\
    Ingénieur d'État en Intelligence Artificielle \\
    ismail.aissaoui.pro@gmail.com \\
    +213 660 70 77 96 \\
    \href{https://ismail-aissaoui.vercel.app}{ismail-aissaoui.vercel.app}
\end{flushleft}

\begin{flushright}
    À l'attention de l'équipe d'encadrement du \\
    \textbf{Laboratoire CESI LINEACT} \\
    Campus de Strasbourg \\
    \medskip
    Le 30 Avril 2026
\end{flushright}

\bigskip

\noindent \textbf{Objet : Candidature pour le doctorat en IA distribuée — Projet FAIR'nRG}

\bigskip

Madame, Monsieur,

Ingénieur d'État en Intelligence Artificielle diplômé de l'ESI-SBA, je vous adresse ma candidature pour le poste de chercheur doctorant sur le sujet "IA distribuée pour la performance énergétique collaborative des prosommateurs". Votre projet FAIR'nRG, au cœur des enjeux de transition énergétique trinationaux, correspond exactement à mon expertise technique et à mes ambitions scientifiques.

Mon parcours m'a permis de développer une double compétence rare, parfaitement alignée avec les axes de votre offre de thèse :

Premièrement, mon travail actuel sur un \textbf{Jumeau Numérique} pour turbines à gaz (Sonatrach) m'a confronté aux défis de l'IA sur périphériques contraints. J'ai implémenté une architecture duale utilisant \textbf{Mamba-SSM} et \textbf{xLSTM} compilé en \textbf{C++ JIT}, atteignant une latence d'inférence de 0.04ms. Cette expertise en optimisation matérielle-logicielle est un atout majeur pour votre premier axe sur l'apprentissage fédéré hétérogène.

Deuxièmement, j'ai une solide expérience dans la conception d'algorithmes d'optimisation complexes. J'ai notamment développé un système d'ordonnancement universitaire basé sur des heuristiques hybrides (\textbf{PSO/ACO/GA}), résolvant des problèmes de contraintes dures à grande échelle. Cette maîtrise des métaheuristiques répond directement à votre second axe sur le pilotage dynamique des systèmes PV+stockage.

Enfin, ma formation en ingénierie logicielle distribuée (Docker, Kubernetes) et ma capacité à travailler sur des projets à forte dimension appliquée me permettront d'être immédiatement opérationnel pour le déploiement du pilote FAIR'nRG.

Intégrer le laboratoire CESI LINEACT et collaborer avec vos partenaires français, allemands et suisses serait pour moi une opportunité unique de contribuer à une recherche souveraine et collaborative. Je suis convaincu que mon autonomie et ma rigueur scientifique sauront répondre à vos attentes.

Je vous remercie par avance de l'attention que vous porterez à mon dossier et je me tiens à votre entière disposition pour un entretien.

Je vous prie d'agréer, Madame, Monsieur, l'expression de mes salutations distinguées.

\bigskip

\begin{flushright}
    \textbf{Ismail AISSAOUI}
\end{flushright}

\end{document}
